% !TeX root = ../main.tex

\chapter{Introduction}\label{ch:introduction}

\section{Clinical Motivation and the Diagnostic Paradigm Shift}\label{sec:motivation}
\outlinenote{Why prediction at the \ac{MCI} stage is the decision point of interest, and
what a conversion label means clinically.}

The primary scope of this thesis is the methodological evaluation of supervised
longitudinal classification approaches for \ac{MCI}-to-\ac{AD} conversion. It does not
claim to develop or validate a normative healthy-reference framework for individualized
abnormality detection.

\section{Resting-State Functional MRI as a Non-Invasive Window}\label{sec:rsfmri}
\outlinenote{Why \ac{FC} is the modality of choice, relative to \ac{CSF} markers,
\ac{PET} and structural \ac{MRI}.}

\section{Deep Learning in Neurodegenerative Disease}\label{sec:dl-neurodegeneration}
\outlinenote{A short survey of what has been built for neurodegenerative prediction from
imaging, paired with the recurring difficulties: small labelled cohorts relative to
feature dimensionality, irregular follow-up, and multi-site heterogeneity. Points forward
to \autoref{ch:foundations} for the detailed treatment.}

\section{Research Questions, Objectives and Contributions}\label{sec:research-questions}
\outlinenote{The four research questions stated once, then the contribution list keyed to
the chapters that deliver each.}

\begin{description}
  \item[RQ1] Does longitudinal modelling add information beyond static \ac{FC}
    representations? Answered in \autoref{sec:rq1}.
  \item[RQ2] Does spatial-first graph representation learning add value?
    Answered in \autoref{sec:rq2}.
  \item[RQ3] Does temporal-first modelling improve over spatial-first modelling?
    Answered in \autoref{sec:rq3}.
  \item[RQ4] Do the findings transfer to an independent cohort?
    Answered in \autoref{sec:rq4}.
\end{description}

\section{Thesis Organisation}\label{sec:organisation}
